\section{Repository audit and diagnosis}\label{sec:audit}

\subsection{The checked logical substrate}

The audited snapshot is commit
\[
\texttt{7c1f2594117c15a980d24300593d625722812768}.
\]
At that snapshot, \RepoPath{Logic/PrimitiveRecursiveArithmetic/ConciseIPRA} provides the ordinary deeply embedded presentation of \iPRA, a native scalar presentation, and kernel-checked translations between them.  The ordinary presentation has:
\begin{itemize}
  \item primitive-recursive function symbols \(\mathsf{PR}(k)\), generated by zero, successor, projections, composition, and primitive recursion;
  \item terms \(\PTm(n)\) over \(n\) free variables;
  \item the theory \(\PRAAx\), consisting of defining equations, successor axioms, and induction for quantifier-free formulas;
  \item intuitionistic natural deduction \(\Prf\) with equality, substitution, weakening, and simultaneous cut.
\end{itemize}
The theorem \texttt{pra\_iff\_native} identifies derivability in this theory with derivability of the translated formula in the native scalar core.  Consequently the most economical development strategy is to construct ordinary \(\Prf\,\PRAAx\) derivations and transfer only at the end.

\subsection{The new derivation library changes the baseline}

The snapshot also contains \RepoPath{Logic/PrimitiveRecursiveArithmetic/Derivations}, a rapidly growing development that was not available to several earlier framework packets.  Its reusable core now derives:
\begin{itemize}
  \item addition, multiplication, predecessor, truncated subtraction, order, parity, powers, finite sums and products, together with their defining equations;
  \item pairing, projections, and simultaneous primitive recursion;
  \item a checked least bounded search operator with object-level bound, validity, and minimality theorems;
  \item division with remainder, gcd, extended Euclid/B\'ezout, bounded-sum reflection, pigeonhole and permutation tools, and generic descent;
  \item a quotient \(Q_\Gamma\) of terms by provable equality, integers as differences, and a tactic \texttt{ipra\_ring} that converts metalevel polynomial normalization into object-level equality derivations.
\end{itemize}
On top of this core, the repository records many actual theorem modules whose public statements have type \(\Prf\,\PRAAx\), including finite algebra, geometry, number theory, and combinatorics.  The exact-algebra and checked-search components are especially important here: they allow \KIFC to concentrate on formula reflection, finite tables, relational composition, and proof scheduling instead of rebuilding a second arithmetic foundation.  The design pattern remains strict: a rich metatheory may discover and normalize a calculation, but the exported theorem must still be a value of the embedded derivability type and its dependency audit must contain no theorem-specific oracle.

\subsection{The proof-packet corpus}

The proof-packet index at the audited snapshot records 78 delivered research packets in algebra, analysis, analytic and Diophantine number theory, geometry, probability, and topology.  Their common strengths are:
\begin{enumerate}[label=(\arabic*)]
  \item they search for finite, quantitative restatements rather than appeal to completed infinitary objects;
  \item they often provide exact-arithmetic reference programs and explicit moduli;
  \item they recognize constructive discontinuities and sometimes replace functions by relations;
  \item they identify theorem-specific missing data, such as height bounds or separation margins.
\end{enumerate}
Their common weaknesses are equally systematic:
\begin{enumerate}[label=(\arabic*)]
  \item many Lean files are specifications or semantic proofs rather than derivations in \(\Prf\,\PRAAx\);
  \item several induction arguments quantify over a bounded range inside the invariant rather than compiling the range to a Boolean state;
  \item coherent approximate witnesses are assumed to follow from pointwise nonemptiness;
  \item exact labels, signs, axes, pivots, or branches are requested at singular inputs;
  \item an ineffective existence theorem is sometimes followed by a bounded enumerator without an effective bound connecting them.
\end{enumerate}

\subsection{Earlier general frameworks}

The foundations packets already contain much of the semantic architecture needed here: certified finitary presentations, finite-use relational composition, demand propagation, proof-carrying observation graphs, stable observables, and finite gluing.  The strongest earlier proposals, proof-producing finitary interface graphs and kernel-integrated contracts, correctly require local object-level certificates and a two-pass compiler.  Their central graph and relational compilers remain blueprints, while the repository now supplies concrete checked algebraic, order, pairing, search, and descent backends.  The present second edition therefore focuses on the missing interfaces and on a dependency-ordered route into the current workspace.

\KIFC therefore does not replace those ideas.  It consolidates them around four stricter rules.

\begin{principle}[Object-proof boundary]
A public theorem is complete only when its conclusion is an embedded derivability judgement.  Computation, model evaluation, and semantic Lean theorems are evidence used to build that judgement, not substitutes for it.
\end{principle}

\begin{principle}[Query specialization]
The theorem and its finite graph are fixed before object-level execution.  The compiler emits dedicated primitive-recursive syntax.  It never asks \iPRA\ to interpret an arbitrary code for an arbitrary primitive-recursive function.
\end{principle}

\begin{principle}[Finite use before infinity]
A theorem is first proved for the finite observations it actually consumes.  A coherent infinite name is extracted only from explicit diameter, motion, and tail data.  Metastable or relational conclusions remain in those forms when no stronger data are present.
\end{principle}

\begin{principle}[Producer/verifier separation]
Deep mathematical input is represented by a finite certificate and a primitive-recursive verifier.  The framework proves the verifier, search, and composition theorems.  It does not rename an unknown quantitative producer as a ``compactness step.''
\end{principle}

\subsection{A failure-mode taxonomy}

For planning purposes, every obstacle in the corpus can be assigned to one of five categories.

\begin{description}[style=nextline,leftmargin=2.2em]
\item[Logical boundary failure.] The claimed result is a metatheoretic proposition rather than \(\Prf\,\PRAAx\).  Remedy: export formulas, local derivations, and an axiom audit.
\item[Induction-shape failure.] The recurrence is finite but the invariant is not quantifier-free.  Remedy: Booleanize bounded content and store enough prefix data in the state.
\item[Representation failure.] The requested operation is discontinuous on the chosen codes.  Remedy: relations, clusters, projective data, stable observers, or an explicit gap.
\item[Effectivity failure.] A theorem asserts existence or finiteness but supplies no primitive-recursive bound, modulus, or constants.  Remedy: state a conditional producer interface and prove the finite verifier path separately.
\item[Coherence failure.] Valid approximants exist independently but do not form a name.  Remedy: backward viability for finite use; shrinking diameter and a summable motion budget for an infinite name.
\end{description}

The rest of the report develops one framework whose types make these failures visible before a large proof is attempted.
